%=========== ΛΥΣΗ ΘΕΜΑΤΟΣ Γ.98 ===========
\begin{thema}{Γ}

\begin{center}
\begin{tikzpicture}[scale=0.7]
  \draw[help lines, gray!30] (-1,-3) grid (7,3);
  \draw[-latex] (-1.3,0) -- (7.3,0) node[right] {$x$};
  \draw[-latex] (0,-3.3) -- (0,3.3) node[above] {$y$};
  \foreach \x in {1,2,3,4,5,6} \draw (\x,0.1) -- (\x,-0.1) node[below,font=\footnotesize] {$\x$};
  \foreach \y in {-2,-1,1,2} \draw (0.1,\y) -- (-0.1,\y) node[left,font=\footnotesize] {$\y$};
  
  \fill[orange!12, opacity=0.5] (0,0) -- (0,2) -- (2,0) -- cycle;
  \fill[red!10, opacity=0.5] (2,0) -- (4,-2) -- (6,0) -- cycle;
  
  \draw[very thick, red!70!black] (0,2) -- (2,0) -- (4,-2) -- (6,0);
  
  \filldraw[red!70!black] (0,2) circle (2pt) node[above right,font=\footnotesize] {$(0,2)$};
  \filldraw[red!70!black] (2,0) circle (2pt) node[above right,font=\footnotesize] {$2$};
  \filldraw[red!70!black] (4,-2) circle (2pt) node[below,font=\footnotesize] {$(4,-2)$};
  \filldraw[red!70!black] (6,0) circle (2pt) node[above right,font=\footnotesize] {$6$};
  \node[red!70!black, fill=white, inner sep=1pt] at (5,2) {$C_{f'}$};
  
  % Area tags
  \node[orange!80!black, font=\footnotesize, fill=white, inner sep=1pt] at (0.8, 0.6) {$E_1 = 2$};
  \node[red!80!black, font=\footnotesize, fill=white, inner sep=1pt] at (4, -0.8) {$E_2 = 4$};
  
  % Monotonicity
  \node[green!50!black, fill=white, inner sep=1pt] at (1, 1.5) {$f'(x) > 0 \ (f \uparrow)$};
  \node[red!80!black, fill=white, inner sep=1pt] at (3.5, -2.5) {$f'(x) < 0 \ (f \downarrow)$};
  \node[green!50!black, fill=white, inner sep=1pt] at (2, 0.5) {\textbf{Τοπικό Μέγιστο}};
\end{tikzpicture}
\end{center}

\begin{erwthma}

\item Από το γράφημα της \textbf{παραγώγου $\mathbf{f'}$} εξετάζουμε το πρόσημό της:
- Στο διάστημα $[0, 2)$, η $f'$ βρίσκεται απολύτως πάνω από τον άξονα των $x$ ($f'(x) > 0$). Άρα η αρχική συνάρτηση $f$ είναι \textbf{γνησίως αύξουσα} στο $[0, 2]$.
- Στο διάστημα $(2, 6]$, η $f'$ βρίσκεται απολύτως κάτω από τον άξονα των $x$ ($f'(x) < 0$, αν και ακουμπά το 0 στο 6). Άρα η $f$ είναι \textbf{γνησίως φθίνουσα} στο $[2, 6]$.
Επειδή η συνάρτηση ξεκινά με συνεχή άνοδο έως το $x=2$ και κατόπιν κατέρχεται χωρίς να ανακάμψει μέχρι το τέλος, στο \textbf{$\mathbf{x_0 = 2}$} παρουσιάζει \textbf{ολικό μέγιστο}.

\item Χρησιμοποιούμε τη σχέση του Θεμελιώδους Θεωρήματος Ολοκληρωτικού Λογισμού:
\[ \int_0^2 f'(x)\,\d x = f(2) - f(0) \]
Γεωμετρικά, το $\int_0^2 f'(x)\,\d x$ είναι το εμβαδόν του αντίστοιχου χωρίου κάτω από το ευθύγραμμο τμήμα της $f'$ (ένα ορθογώνιο τρίγωνο).
- Βάση = 2 ($0$ έως $2$).
- Ύψος = $2$ ($f'(0) = 2$).
Άρα $\text{Εμβαδόν} = \frac{1}{2} \cdot 2 \cdot 2 = 2$.
Συνεπώς:
\[ f(2) - f(0) = 2 \xRightarrow{f(0)=0} \mathbf{f(2) = 2} \]

\item Ζητείται η εξίσωση της εφαπτομένης της $C_f$ στο $x_0 = 0$.
Μας δίνονται αμέσως τα απαραίτητα στοιχεία:
- Σημείο επαφής: $(0, f(0)) = (0, 0)$.
- Κλίση (από το γράφημα της $f'$): $\lambda = f'(0) = 2$.
Εφαρμόζοντας τον τύπο της εφαπτομένης:
\[ y - f(0) = f'(0) (x - 0) \Rightarrow y - 0 = 2(x - 0) \Rightarrow \mathbf{y = 2x} \]

\item Μας δίνεται ο τύπος υπαναχώρησης αξίας:
\[ f(6) = f(2) + \int_2^6 f'(x)\,\d x \]
Γνωρίζουμε ήδη ότι $f(2) = 2$. Υπολογίζουμε ξανά γεωμετρικά το εμβαδόν του τριγώνου για το ολοκλήρωμα $\int_2^6$:
- Βάση = $6 - 2 = 4$.
- Ύψος (κάτω από τον άξονα) = $2$ (φτάνει στο -2, άρα το μήκος είναι 2).
Επειδή το τρίγωνο βρίσκεται κάτω από τον οριζόντιο άξονα, το προσημασμένο ολοκλήρωμα είναι αρνητικό:
\[ \int_2^6 f'(x)\,\d x = - \left( \frac{1}{2} \cdot \text{βάση} \cdot \text{ύψος} \right) = - \left( \frac{1}{2} \cdot 4 \cdot 2 \right) = -4 \]
Αντικαθιστούμε:
\[ f(6) = 2 - 4 = \mathbf{-2} \]
Πράγματι, διαπιστώνουμε ότι \textbf{$\mathbf{f(6) < 0}$} (συγκεκριμένα $-2 < 0$), αφού η αρνητική επιφάνεια του δεύτερου τριγώνου (-4) "υπερκέρασε" τη θετική (2), ωθώντας την αρχική συνάρτηση $f$ πιο κάτω από το 0.
\end{erwthma}
\end{thema}
